% TEMPLATE-MILC-v3.tex - plantilla maestra UNICA de las guias MILC v3.
%
% La usan las cuatro familias de guias, cada una con su driver:
%   · Grados 6-11          -> scripts/build-guias-g11.py
%   · Bebras               -> scripts/build-guias-bebras.py
%   · Semillero            -> scripts/build-guias-semillero.py
%   · Territorio interior  -> scripts/build-guias-territorio-interior.py
%
% Antes habia tres copias casi identicas (~400 lineas iguales, 6 distintas):
% cada mejora habia que aplicarla tres veces, y en la practica no se hacia
% (el motor de recursos vivia solo en dos de ellas). Lo que varia entre
% familias esta parametrizado en 6 placeholders que cada driver llena
% (se nombran aqui SIN su sintaxis real: el builder reemplaza por texto plano,
% tambien dentro de los comentarios, y un valor multilinea rompe el preambulo):
%
%   HEADER_IZQ        encabezado izquierdo de pagina
%   HEADER_DER        encabezado derecho de pagina
%   PORTADA_ETIQUETA  rotulo sobre el numero grande (GRADO/MOMENTO/MODULO)
%   PORTADA_SERIE     linea de serie de la portada
%   PORTADA_ID        identificador de la guia en la portada
%   PORTADA_CIRCULO   el circulito de grado (°), o vacio si no aplica
%   PDF_TITULO        titulo en texto plano (sin \textbf ni comillas TeX) para
%                     los metadatos del PDF (/Title); lo produce lib_xelatex.texto_pdf
%
% Composicion (auditoria 2026-09-03): las bandas de estacion piden espacio
% (\Needspace*) y prohiben el salto justo despues (\nopagebreak), asi no quedan
% huerfanas al pie; las cajas largas (softbox, drawbox, infoband, puentebox) son
% `breakable`, asi una caja de media pagina ya no arrastra todo a la siguiente
% dejando la anterior al 40 %. Todo texto sobre mostaza o verde va en milcNegro
% (blanco sobre #E5B400 da 1,93:1 y sobre #58A923 2,95:1; ninguno pasa WCAG AA).
% El resto de claves las documenta content/guias/_SCHEMA.md. El script valida
% que no queden placeholders sin reemplazar antes de escribir el archivo final.

% Etiquetado PDF (latex-lab): /Lang, estructura y texto alternativo de las
% figuras, para lectores de pantalla. Debe ir ANTES de \documentclass.
\DocumentMetadata{lang=es-CO,tagging=on}
\documentclass[11pt,letterpaper]{article}
% headheight=24pt: la cabecera derecha lleva dos lineas (identidad de la guia
% y estacion actual). headsep se acorta en la misma medida para que el bloque
% de texto quede exactamente donde estaba con headheight=12pt.
\usepackage[margin=2.0cm,headheight=24pt,headsep=13pt]{geometry}
\usepackage[table]{xcolor}
\usepackage{fontspec}
\usepackage[spanish]{babel}
\usepackage{tikz}
% Librerías TikZ para los diagramas de las guías (bloque `recursos:`):
%  · babel  → babel en español vuelve activos caracteres como «>», y TikZ los usa
%             en sus opciones (>=stealth, ->). Sin esta librería, cualquier
%             diagrama con flechas revienta con «\language@active@arg> has an
%             extra }». Debe ir ANTES de que se dibuje nada.
%  · positioning        → sintaxis «below left=2pt and 3pt».
%  · decorations.pathreplacing → llaves ({brace}) para señalar rangos.
\usetikzlibrary{babel,positioning,decorations.pathreplacing}
\usepackage{graphicx}
% Circuitos: el trabajo de electrónica se hace hoy con micro:bit y sus sensores
% integrados (MakeCode, sin cableado), así que no se necesitan esquemáticos.
% Si algún día entran montajes con protoboard, basta descomentar esta línea:
% los diagramas .tex/.tikz declarados en `recursos:` ya se incrustan solos.
% \usepackage{circuitikz}
\usepackage[most]{tcolorbox}
\usepackage{tabularx}
\usepackage{array}
\usepackage{fancyhdr}
\usepackage{titlesec}
\usepackage[document]{ragged2e}  % bandera a la derecha en todo el cuerpo
\usepackage{enumitem}
\usepackage{microtype}
\usepackage{etoolbox}
\usepackage{needspace}
% hyperref al final del preambulo: metadatos del PDF (titulo, autor, idioma,
% familia), marcadores por estacion (\pdfbookmark) y enlaces sin recuadro.
\usepackage[hidelinks,
  pdftitle={Gráficos que no engañan — barras, líneas y circular con el eje en cero},
  pdfauthor={PhD. Álvaro Cárdenas Orozco},
  pdfsubject={Tecnología e Informática · Grado 8 · Guía 7 · MILC v3},
  pdflang={es-CO},
  bookmarksopen=true,bookmarksnumbered=false]{hyperref}

% Helvetica Neue no trae la flecha U+2192, asi que cada «→» escrito literalmente
% en el YAML se compilaba a NADA, en silencio (462 flechas en 61 guias: rutas del
% tipo «paleta → tipografia → lockup» salian rotas). Mapeamos el caracter a su
% version matematica, que si tiene glifo. Asi el autor puede seguir escribiendo
% «→» comodamente en el YAML.
\usepackage{newunicodechar}
\newunicodechar{→}{\ensuremath{\rightarrow}}
% Mismo problema con los simbolos de marca y vinetas: el «✓» de «una palomita ✓»
% salia como cuadrito vacio en el PDF de todas las guias que lo usaban.
\usepackage{pifont}
\newunicodechar{✓}{\ding{51}}
\newunicodechar{✗}{\ding{55}}
\newunicodechar{✔}{\ding{52}}
\newunicodechar{•}{\textbullet}
\newunicodechar{≥}{\ensuremath{\geq}}
\newunicodechar{≤}{\ensuremath{\leq}}

\setmainfont{Helvetica Neue}
\setsansfont{Helvetica Neue}
\setlength{\parindent}{0pt}
\setlength{\parskip}{6pt}
% Sin particion de palabras: en 6.º-7.º una palabra cortada por guion al final
% de linea («Comuni-cacion») frena la lectura mucho mas que una linea corta.
\hyphenpenalty=10000
\exhyphenpenalty=10000
\pretolerance=2000
\tolerance=3000
\setlength{\emergencystretch}{3em}
\linespread{1.10}
\renewcommand{\arraystretch}{1.25}
\setlist{leftmargin=5mm,itemsep=1mm,topsep=1mm}
\newcolumntype{Y}{>{\RaggedRight\arraybackslash}X}

\definecolor{milcMagenta}{HTML}{E6007E}
\definecolor{milcVino}{HTML}{5A0038}
\definecolor{milcTurquesa}{HTML}{0093A5}
\definecolor{milcVerde}{HTML}{58A923}
\definecolor{milcMostaza}{HTML}{E5B400}
\definecolor{milcOcre}{HTML}{B86600}
\definecolor{milcNegro}{HTML}{241816}
\definecolor{milcGris}{HTML}{F7F7F4}
\definecolor{milcLinea}{HTML}{E8E2DD}
\definecolor{milcGradePrimary}{HTML}{FF6600}
\definecolor{milcGradeDark}{HTML}{9A3E00}
\definecolor{milcGradeSoft}{HTML}{FFE4D0}
\definecolor{milcGradeAccent}{HTML}{CFE7FF}
\definecolor{milcGradeText}{HTML}{FFFFFF}
\color{milcNegro}

\pagestyle{fancy}
\fancyhf{}
% Cabecera de dos lineas: identidad de la guia arriba y, a la derecha y en
% gris, la estacion en curso (\markright desde \milcestacion). Antes de la
% primera estacion la segunda linea va vacia; el \strut mantiene la altura
% para que la linea de arriba no baile entre paginas.
\lhead{\small Tecnología e Informática · Grado 8\\[-1pt]{\footnotesize\strut}}
\rhead{\small Guía 7 · MILC v3\\[-1pt]{\footnotesize\strut\color{milcNegro!65}\rightmark}}
\cfoot{\small\thepage}
\renewcommand{\headrulewidth}{0pt}

% Sobre mostaza (#E5B400) y verde (#58A923) el texto va en milcNegro; sobre el
% resto de la paleta, en blanco. Se usa dentro de fonttitle/fontupper, donde un
% \color posterior manda sobre coltitle.
\newcommand{\milctintasobre}[1]{%
  \ifboolexpr{test{\ifstrequal{#1}{milcMostaza}} or test{\ifstrequal{#1}{milcVerde}}}%
    {\color{milcNegro}}{\color{white}}}

\newtcolorbox{titlebox}[1]{
  enhanced,colback=#1,colframe=#1,arc=7mm,boxrule=0pt,
  left=7mm,right=7mm,top=3mm,bottom=3mm,
  before skip=8pt,after skip=8pt,
  fontupper=\sffamily\bfseries\Large\color{white}
}

\newtcolorbox{softbox}[3]{
  enhanced,breakable,pad at break=3mm,
  colback=#2,colframe=#1,borderline west={4pt}{0pt}{#1},
  arc=2mm,boxrule=.4pt,left=4mm,right=4mm,top=3mm,bottom=3mm,
  before skip=6pt,after skip=8pt,title={#3},coltitle=white,
  colbacktitle=#1,fonttitle=\sffamily\bfseries\milctintasobre{#1},fontupper=\RaggedRight,
  attach boxed title to top left={xshift=3mm,yshift=-2mm},
  boxed title style={arc=2mm, boxrule=0pt}
}

\newtcolorbox{drawbox}[2]{
  enhanced,breakable,pad at break=4mm,
  colback=white,colframe=#1,arc=4mm,boxrule=1pt,
  left=5mm,right=5mm,top=4mm,bottom=4mm,before skip=8pt,after skip=8pt,
  title={#2},coltitle=white,colbacktitle=#1,fonttitle=\sffamily\bfseries\milctintasobre{#1},
  fontupper=\RaggedRight,
  attach boxed title to top left={xshift=4mm,yshift=-2mm},
  boxed title style={arc=3mm, boxrule=0pt}
}

\newtcolorbox{infoband}[1]{
  enhanced,breakable,pad at break=2mm,colback=#1!8,colframe=#1!50,arc=2mm,boxrule=.5pt,
  left=4mm,right=4mm,top=2mm,bottom=2mm,before skip=4pt,after skip=4pt,
  fontupper=\small\RaggedRight
}

% Puente narrativo entre secciones (contrato editorial Conéctate).
% Da continuidad: "Acabas de X. Ahora vas a Y porque Z."
% Visualmente: banda izquierda en color del grado, texto en itálica pequeña.
\newtcolorbox{puentebox}{
  enhanced,breakable,pad at break=2mm,colback=milcGradeSoft,colframe=milcGradeDark,
  borderline west={3pt}{0pt}{milcGradeDark},
  arc=0mm,boxrule=0pt,left=5mm,right=4mm,top=2.5mm,bottom=2.5mm,
  before skip=4pt,after skip=8pt,
  fontupper=\itshape\small\color{milcNegro}\RaggedRight
}
% La flecha va en modo matematico a proposito: Helvetica Neue NO tiene el glifo
% U+2192, asi que \textrightarrow se compilaba a NADA («Missing character: There
% is no → in font Helvetica Neue») y el puente salia sin flecha. En modo
% matematico la flecha viene de la fuente matematica y si se dibuja.
\newcommand{\puente}[1]{%
  \begin{puentebox}%
    \textcolor{milcGradeDark}{$\rightarrow$}\hspace{2mm}#1%
  \end{puentebox}%
}

\newcommand{\linead}[1]{\noindent\color{milcLinea}\rule{#1}{0.5pt}\color{milcNegro}}
\newcommand{\respuesta}[1]{\par\vspace{2mm}\foreach \n in {1,...,#1}{\noindent\color{milcLinea}\rule{\linewidth}{0.5pt}\par\vspace{5mm}}\color{milcNegro}}
\newcommand{\checkbox}{\textcolor{milcGradeDark}{\fbox{\phantom{x}}}\hspace{5pt}}

% ─── Listas aireadas para las guías socioemocionales ────────────────────
% Componer las verificaciones como «\checkbox texto\\» deja el texto que salta
% de línea debajo del cuadrito y apelmaza el bloque. Estos entornos alinean la
% continuación y airean los ítems. Los usa Territorio Interior; cualquier
% familia puede hacerlo.
\newlist{checklista}{itemize}{1}
\setlist[checklista]{
  label=\textcolor{milcGradeDark}{\fbox{\phantom{x}}},
  leftmargin=9mm, labelsep=3.2mm, itemsep=2.4mm,
  topsep=2.5mm, parsep=0pt, partopsep=0pt, before=\RaggedRight
}
\newlist{pasoslista}{enumerate}{1}
\setlist[pasoslista]{
  label=\textcolor{milcGradeDark}{\sffamily\bfseries\arabic*.},
  leftmargin=8mm, labelsep=3mm, itemsep=2.8mm,
  topsep=1mm, parsep=0pt, partopsep=0pt, before=\RaggedRight
}

\newcommand{\phasecard}[4]{
\begin{tcolorbox}[enhanced,colback=#3,colframe=#2,arc=3mm,boxrule=.5pt,height=3.05cm,valign=center,halign=center,left=1.5mm,right=1.5mm,top=2mm,bottom=2mm]
{\sffamily\bfseries\fontsize{22}{24}\selectfont\textcolor{#2}{#1}}\par
{\sffamily\bfseries\small #4}
\end{tcolorbox}
}

% ── Piezas del rediseno 2026 (legibilidad 6.º-11.º) ───────────────────
% Banda de estacion: reemplaza el titlebox macizo. Titulo a la izquierda,
% dato practico (tiempo/modalidad) a la derecha, en una sola linea.
% Nunca huerfana: pide 9 lineas libres (o salta de pagina rellenando con
% \vfill) y prohibe el salto justo despues de la banda. Nueve y no seis:
% la caja partible que sigue necesita ~80pt (titulo adosado, relleno y dos
% lineas) para arrancar en la misma pagina; con menos, tcolorbox la manda
% entera a la siguiente y la banda se queda sola. El penalty 10000 va
% en `after`, ANTES del espacio posterior: un `after skip` (aunque sea 0pt)
% es un \vskip tras la caja, y ese glue es un punto de corte legal que un
% \nopagebreak escrito despues de \end{tcolorbox} ya no cierra. Cada banda
% deja marcador PDF y actualiza la cabecera derecha.
\newcounter{milcestacion}
\newcommand{\milcestacion}[2]{%
\Needspace*{9\baselineskip}%
\stepcounter{milcestacion}%
\pdfbookmark[1]{#2}{milcestacion\themilcestacion}%
\markright{#2}%
\begin{tcolorbox}[enhanced,colback=#1,colframe=#1,arc=2mm,boxrule=0pt,
  left=5mm,right=5mm,top=2.6mm,bottom=2.6mm,before skip=11pt,
  after={\par\nopagebreak\vskip7pt}]
{\sffamily\bfseries\fontsize{15}{18}\selectfont\milctintasobre{#1}#2}
\end{tcolorbox}}

% Tarjeta de paso numerada: sustituye las tablas «Pilar / Aplicacion», que
% eran formato de consulta docente, no de estudiante. El numero va en un
% circulo sobre el borde para que la secuencia se lea de un vistazo.
\newcommand{\milcpaso}[3]{%
\Needspace{4\baselineskip}%
\begin{tcolorbox}[enhanced,colback=white,colframe=milcLinea,arc=1.5mm,boxrule=.9pt,
  left=14mm,right=4mm,top=3mm,bottom=3mm,before skip=4pt,after skip=4pt,
  fontupper=\RaggedRight,
  overlay={\node[anchor=north west,circle,fill=#2,text=white,inner sep=0pt,
    minimum size=7.4mm,font=\sffamily\bfseries\small]
    at ([xshift=3.6mm,yshift=-3.2mm]frame.north west) {#1};}]
#3
\end{tcolorbox}}

% Casilla marcable de tamano util con lapiz (la anterior era \fbox{\phantom{x}},
% de alto variable segun el texto que la seguia).
\newcommand{\milcmarca}{\framebox[3.8mm]{\rule{0pt}{3.4mm}}}

% Fila marcable de lista suelta (checks de la estacion 1).
\newcommand{\milccheckrow}[1]{%
\par\vspace{1.8mm}\noindent
\makebox[6.4mm][l]{\raisebox{-.55mm}{\textcolor{milcNegro}{\milcmarca}}}%
\parbox[t]{\dimexpr\linewidth-6.4mm\relax}{\RaggedRight #1}\par}

% Celda marcable dentro de la rubrica (tabla de autoevaluacion). Misma
% sangria francesa, pero pensada para vivir dentro de una columna X.
\newcommand{\milcopcion}[1]{%
\makebox[5.2mm][l]{\raisebox{-.55mm}{\textcolor{milcNegro}{\milcmarca}}}%
\parbox[t]{\dimexpr\linewidth-5.2mm\relax}{\RaggedRight #1}}

% ── Recursos visuales de la guía ──────────────────────────────────────
% Declarados en el bloque `recursos:` del YAML y emitidos por el builder.
% \guiaFigura   → imagen rasterizada o PDF (foto, carta celeste, montaje).
% \guiaDiagrama → fuente TikZ/circuitikz (.tex) incrustada como vector:
%                 es la vía para circuitos y esquemáticos editables.
% [#1] = texto alternativo (alt) · #2 = ruta relativa al .tex · #3 = pie
% El alt llega al PDF etiquetado (/Alt) y lo lee el lector de pantalla.
\newcommand{\guiaFigura}[3][]{%
\begin{center}
\includegraphics[width=0.88\linewidth,height=0.38\textheight,keepaspectratio,alt={#1}]{#2}\\[1.5mm]
{\footnotesize\itshape\textcolor{milcNegro!75}{#3}}
\end{center}
}

\newcommand{\guiaDiagrama}[2]{%
\begin{center}
\input{#1}\\[1.5mm]
{\footnotesize\itshape\textcolor{milcNegro!75}{#2}}
\end{center}
}

\begin{document}
\thispagestyle{empty}

% ── Portada ───────────────────────────────────────────────────────────
% Antes: un rectangulo de color a pagina completa. Se veia bien en pantalla,
% pero en la fotocopiadora del colegio se comia el toner y salia gris sucio.
% Ahora la banda ocupa el tercio superior y el resto va en blanco.
\begin{tikzpicture}[remember picture,overlay]
  \path[fill=milcGradePrimary]
    (current page.north west) rectangle ([yshift=-10.6cm]current page.north east);

  \node[anchor=north west,text=milcGradeText] at ([xshift=2.0cm,yshift=-1.55cm]current page.north west)
    {\sffamily\bfseries\fontsize{12}{14}\selectfont GRADO};
  \node[anchor=north west,text=milcGradeText,opacity=.85] at ([xshift=2.0cm,yshift=-2.35cm]current page.north west)
    {\sffamily\bfseries\fontsize{8.5}{10}\selectfont SERIE GUÍAS · TECNOLOGÍA E INFORMÁTICA};
  \node[anchor=north east,text=milcGradeText,opacity=.85,align=right] at ([xshift=-2.0cm,yshift=-1.55cm]current page.north east)
    {\sffamily\bfseries\fontsize{9}{12}\selectfont I.E. Sor María Juliana\\Cartago, Valle del Cauca};

  \node[anchor=west,text=milcGradeText] (gradonum) at ([xshift=1.85cm,yshift=-7.1cm]current page.north west)
    {\sffamily\bfseries\fontsize{112}{112}\selectfont 8};
  % El circulito de grado (°) solo aplica a las guias de grado («11°»); en
  % Bebras y Semillero el numero grande es un momento/modulo, no un grado.
  % Se posiciona relativo al nodo `gradonum`, no en coordenadas absolutas.
  \draw[milcGradeText,line width=1.6pt]
    ([xshift=2.0mm,yshift=-4.5mm]gradonum.north east) circle (.15cm);

  \node[anchor=west,text=milcGradeText,text width=11.4cm,align=left] at ([xshift=1.5cm,yshift=0.5cm]gradonum.east)
    {
     {\sffamily\bfseries\fontsize{9}{11}\selectfont\MakeUppercase{Período 1 · Análisis de datos con phronesis}}\\[3mm]
     {\sffamily\bfseries\fontsize{21}{25}\selectfont\setlength{\baselineskip}{27pt}Gráficos que\\[4pt]no engañan}\\[3.5mm]
     {\sffamily\bfseries\fontsize{15}{18}\selectfont Guía 7-8-TIC}
    };
\end{tikzpicture}

\vspace*{9.4cm}
\vfill

{\sffamily\bfseries\fontsize{8}{10}\selectfont\textcolor{milcGradeDark}{LO QUE TE LLEVAS HOY}}

\begin{tcolorbox}[enhanced,colback=white,colframe=milcNegro,arc=2mm,boxrule=1.6pt,
  left=5mm,right=5mm,top=4mm,bottom=4mm,before skip=3pt,after skip=12pt,
  fontupper=\RaggedRight\fontsize{11}{15.5}\selectfont]
Una hoja de Excel con tres gráficos del mismo conjunto de datos, de barras, de líneas y circular, cada uno con título, ejes con unidades y el eje vertical desde cero. Un párrafo por gráfico que diga qué pregunta responde mejor. Y la crítica de un gráfico engañoso real, de prensa o de redes, con la trampa identificada y la corrección propuesta.
\end{tcolorbox}

\begin{tcolorbox}[enhanced,colback=milcGris,colframe=milcGris,arc=2mm,boxrule=0pt,
  left=5mm,right=5mm,top=3.5mm,bottom=3.5mm,after skip=12pt]
{\sffamily\bfseries\fontsize{8}{10}\selectfont\textcolor{milcNegro!55}{ESTA GUÍA ES DE}}\\[3mm]
\begin{tabularx}{\linewidth}{X p{3.2cm} p{3.2cm}}
\linead{\linewidth} & \linead{3.0cm} & \linead{3.0cm}\\[-1mm]
{\fontsize{8.5}{10}\selectfont\textcolor{milcNegro!55}{Nombre completo}} &
{\fontsize{8.5}{10}\selectfont\textcolor{milcNegro!55}{Grupo}} &
{\fontsize{8.5}{10}\selectfont\textcolor{milcNegro!55}{Fecha}}\\
\end{tabularx}
\end{tcolorbox}

\vfill

\noindent\textcolor{milcLinea}{\rule{\linewidth}{1pt}}\\[2mm]
{\sffamily\bfseries\fontsize{9.5}{12}\selectfont Autor: Dr. Álvaro Cárdenas Orozco}\\[1mm]
{\sffamily\fontsize{8.5}{11}\selectfont\textcolor{milcNegro!55}{Metodología MILC · apertura ancestral · pensamiento computacional · triángulo Dussel–estoicismo–Floridi}}

\newpage

\section*{Apertura: Gráficos que no engañan --- barras, líneas y circular con el eje en cero}

\begin{softbox}{milcGradeDark}{milcGradeSoft}{Saber ancestral}
Riosucio, en Caldas, nació de dos pueblos que no se querían. Quiebralomo, de mineros blancos, negros y mulatos, y La Montaña, de indígenas. Cada uno puso su plaza y su iglesia a una cuadra del otro. Entre las dos levantaron una cerca que los separó durante más de veinte años; los días de mercado se insultaban de lado a lado (Ministerio de Cultura, 2011). En 1846 un decreto los fusionó, y de aquellos insultos salió el verso burlón que hoy es el Carnaval. Fíjate en la cerca. La misma ciudad, vista desde una plaza, era un pueblo; vista desde la otra, era otro. Un gráfico hace lo mismo con los datos. Los mismos números, con el eje empezando en cero, cuentan una historia; con el eje empezando en 50, cuentan otra. La \textbf{cara de exclusión}: el relato oficial del Eje Cafetero se contó como el de colonos blancos y respetables. Ese relato dejó por fuera a los indígenas y a los negros que ya estaban ahí (Appelbaum, 2007). Hoy vas a hacer gráficos que no levanten cercas: que dejen ver los datos como son.\par\smallskip{\footnotesize\textcolor{milcNegro!70}{\textbf{Fuente.} Ministerio de Cultura de Colombia. (2011). \emph{Plan Especial de Salvaguardia del Carnaval de Riosucio}, pp. 12--13, 20 y 41.}}
\end{softbox}

\begin{softbox}{milcTurquesa}{milcGris}{Saber contemporáneo}
Un gráfico es una forma de contar los datos con la vista. Las \textbf{barras} comparan valores entre categorías: ventas por día, notas por estudiante. Las \textbf{líneas} muestran cómo cambia algo en el tiempo: temperatura por día, ventas por mes. El \textbf{circular} muestra qué parte del total es cada cosa, y solo sirve con pocas categorías. Cada tipo responde una pregunta distinta. Pero incluso el tipo correcto puede engañar. Un eje vertical que empieza en 50 en vez de cero hace que una diferencia de dos puntos parezca enorme. Un efecto 3D cambia las proporciones. Un circular con diez tajadas no se lee. Un gráfico sin título ni unidades se interpreta como cada quien quiera. Un gráfico honesto tiene cinco cosas: el tipo según la pregunta, el eje desde cero, título que diga qué muestra, ejes con unidades y la fuente de los datos.
\end{softbox}

\begin{softbox}{milcMostaza}{milcGris}{Pregunta-puente}
\textit{Desde cada plaza, Riosucio parecía otro pueblo. Cuando hagas un gráfico de las notas del salón con el eje empezando en 3,0, ¿qué pueblo le estás mostrando a quien lo mira? ¿Y cómo se vería con el eje en cero?}
\end{softbox}

\begin{softbox}{milcVerde}{milcGris}{Saber y saber hacer}
Al terminar podrás: (1) \textbf{identificar} qué tipo de gráfico responde cada pregunta; (2) \textbf{crear} tres gráficos honestos en Excel sobre el mismo conjunto; (3) \textbf{analizar} cuál de los tres comunica mejor y por qué; (4) \textbf{evaluar} un gráfico real de prensa o de redes y decir qué trampa usa.
\end{softbox}



\small
\begin{tabularx}{\textwidth}{>{\bfseries}p{3.0cm}Y}
\rowcolor{milcGradePrimary!12}Autor & Dr. Álvaro Cárdenas Orozco\\
DBA & Tratamiento de datos cuantitativos con criterio ético (MEN, T\&I)\\
Referentes & Phronesis aristotélica · Floridi (ética del dato) · Estoicismo (ver lo que es)\\
Producto final & Una hoja de Excel con tres gráficos del mismo conjunto de datos, de barras, de líneas y circular, cada uno con título, ejes con unidades y el eje vertical desde cero. Un párrafo por gráfico que diga qué pregunta responde mejor. Y la crítica de un gráfico engañoso real, de prensa o de redes, con la trampa identificada y la corrección propuesta.\\
\end{tabularx}
\normalsize

\puente{En la sesión 6 escribiste fórmulas que se leen sin adivinar el orden. Hoy conviertes cifras en gráficos que se leen sin engaño. En la sesión 8 vas a hacer que la hoja avise sola cuando un valor se sale de rango.}

\milcestacion{milcGradeDark}{Ruta MILC: así vamos a aprender}
\begin{tabularx}{\textwidth}{>{\centering\arraybackslash}X>{\centering\arraybackslash}X>{\centering\arraybackslash}X>{\centering\arraybackslash}X}
\phasecard{1}{milcVerde}{milcGris}{Escucha\\[-1mm]\small Observo, escucho y pregunto.} &
\phasecard{2}{milcTurquesa}{milcGris}{Entender\\[-1mm]\small Sistematizo y ordeno.} &
\phasecard{3}{milcMagenta}{milcGris}{Hacer\\[-1mm]\small Construyo y aplico.} &
\phasecard{4}{milcVino}{milcGris}{Cierre\\[-1mm]\small Evalúo y reflexiono.}
\end{tabularx}


\puente{Antes de hacer gráficos, vas a mirar tres que te entrega tu docente y a decir cuál entiendes, cuál exagera y cuál confunde. Solo mirar.}

\milcestacion{milcVerde}{Estación 1 de 3 · Escucha}

\begin{softbox}{milcVerde}{milcGris}{Escena para escuchar}
\textbf{Actividad 1 · IDENTIFICA --- Tres gráficos, un solo dato} (15 min · individual). (1) Mira los tres gráficos que proyecta tu docente: son los mismos datos, dibujados de tres maneras. Uno de barras con el eje en cero, uno de barras con el eje empezando alto, y uno circular en 3D con ocho tajadas. (2) Escribe cuál entiendes mejor a primera vista. (3) Escribe cuál parece exagerar las diferencias y por qué. (4) Escribe cuál confunde más que aclara. (5) Anota una regla que sacarías de lo que viste.
\end{softbox}

{\sffamily\bfseries\fontsize{8}{10}\selectfont\textcolor{milcNegro!55}{MARCA CUANDO LO HAYAS HECHO}}
\milccheckrow{Escribí cuál gráfico entiendo mejor y por qué.}
\milccheckrow{Encontré el que exagera y dije qué lo hace exagerar.}
\milccheckrow{Saqué una regla propia de lo que vi.}

\begin{infoband}{milcVerde}


\textbf{Lo que va al cuaderno (Actividad 1 · IDENTIFICA).} \textbf{Título:} «Actividad 1 --- Tres gráficos, un solo dato». \textbf{Formato:} tabla de 3 filas y 2 columnas (gráfico / qué me hace ver o creer), más una regla propia. \textbf{Extensión:} un tercio de página. \textbf{Sabes que terminaste cuando} las tres filas están llenas y tu regla dice algo que un compañero podría aplicar.
\end{infoband}

\puente{Ya viste que el mismo dato se puede mostrar de tres maneras. Ahora aprendes los tres tipos, las cinco reglas de un gráfico honesto, y con tu pareja construyes los tres sobre tu tabla.}

\milcestacion{milcTurquesa}{Estación 2 de 3 · Entender}

\begin{softbox}{milcTurquesa}{milcGris}{Lo que vas a entender}
\textbf{Actividad 2 · CREA --- Tres gráficos honestos sobre tu tabla} (30 min · parejas). (1) Con tu pareja, abran la tabla limpia de la sesión 3. Elijan una columna numérica con su categoría, por ejemplo ventas por día o notas por estudiante. (2) Seleccionen el rango, Insertar, Gráfico, y hagan uno de barras. (3) Con los mismos datos hagan uno de líneas y uno circular. (4) A cada uno pónganle título que diga qué muestra, ejes con unidades y eje vertical desde cero. (5) Escriban un párrafo por gráfico: qué pregunta responde mejor y por qué.
\end{softbox}

\milcpaso{1}{milcTurquesa}{\textbf{Los tres tipos.} Barras para comparar categorías: ¿quién vendió más? Líneas para ver cambios en el tiempo: ¿cómo cambió de lunes a viernes? Circular para partes de un total: ¿qué parte de la plata se fue en tienda? Con más de seis tajadas, el circular no se lee.}
\milcpaso{2}{milcTurquesa}{\textbf{Las seis trampas.} Eje vertical que no empieza en cero. Escalas distintas en gráficos que se comparan. Efecto 3D que cambia las proporciones. Demasiadas categorías. Colores que sugieren un orden que no existe. Sin título, sin unidades o sin fuente.}
\milcpaso{3}{milcTurquesa}{\textbf{Las cinco reglas del gráfico honesto.} (1) Tipo según la pregunta. (2) Eje vertical desde cero en las barras. (3) Título que diga qué muestra. (4) Ejes con unidades. (5) Fuente de los datos. Si una falta, el gráfico deja de ser un dato y se vuelve una opinión.}
\milcpaso{4}{milcTurquesa}{\textbf{Cómo se hace en Excel.} Selecciona el rango con su encabezado. Insertar, Gráficos, y elige el tipo. Clic derecho en el eje vertical, Dar formato al eje, y pon el mínimo en cero. Doble clic en el título para escribirlo. Léelo en voz alta: si necesitas explicarlo, le falta algo.}

\begin{drawbox}{milcTurquesa}{El gráfico honesto, en una mirada}
\textbf{Tipo:} según la pregunta. \\[1mm] \textbf{Eje vertical:} desde cero. \\[1mm] \textbf{Título:} qué muestra. \\[1mm] \textbf{Ejes:} con unidades. \\[1mm] \textbf{Fuente:} de dónde salen los datos.
\end{drawbox}

\begin{softbox}{milcMostaza}{milcGris}{Errores comunes y su corrección}
(1) \textbf{Eje que empieza alto}: una diferencia de dos puntos parece de veinte. (2) \textbf{3D}: las tajadas de adelante se ven más grandes de lo que son. (3) \textbf{Circular con muchas categorías}: nadie distingue las tajadas. (4) \textbf{Líneas para datos sin orden en el tiempo}: la línea sugiere una tendencia que no existe. (5) \textbf{Sin título ni unidades}: cada quien lee lo que quiere.
\end{softbox}

\begin{infoband}{milcTurquesa}


\textbf{Lo que va al cuaderno (Actividad 2 · CREA).} \textbf{Título:} «Actividad 2 --- Tres gráficos honestos». \textbf{Formato:} tres párrafos de tres líneas, uno por gráfico, con la pregunta que responde mejor; y una marca de cuál de los tres eligieron como el mejor para esos datos. \textbf{Extensión:} media página. \textbf{Sabes que terminaste cuando} los tres gráficos cumplen las cinco reglas y los tres párrafos nombran una pregunta distinta.
\end{infoband}

\puente{Ya sabes construir. Toca criticar: buscas un gráfico engañoso de verdad, nombras la trampa y propones cómo arreglarlo. Y revisas los gráficos de tu pareja con las cinco reglas.}

\milcestacion{milcMagenta}{Estación 3 de 3 · Hacer}

\begin{softbox}{milcMagenta}{milcGris}{Cómo vas a construir tu producto}
\textbf{Actividad 3 · EVALÚA --- Un gráfico engañoso de verdad} (25 min · individual). (1) Busca un gráfico en un periódico, una revista, una publicidad o una red social. Si no hay conexión, tu docente trae varios. (2) Pásalo por las cinco reglas y marca cuáles cumple y cuáles no. (3) Nombra la trampa: cuál de las seis usa. (4) Escribe qué impresión falsa deja en quien lo mira. (5) Propón la corrección: qué cambio pequeño lo haría honesto. (6) Revisa los tres gráficos de tu pareja con las cinco reglas y anota lo que falte.
\end{softbox}

\milcpaso{1}{milcMagenta}{\textbf{Tu entrega tiene tres piezas.} (1) Los tres gráficos con sus párrafos. (2) La crítica del gráfico engañoso: trampa, impresión falsa y corrección. (3) La revisión de los gráficos de tu pareja con las cinco reglas.}
\milcpaso{2}{milcMagenta}{\textbf{Cómo se critica un gráfico.} Primero se describe: qué muestra, qué tipo es, dónde empieza el eje. Después se nombra la trampa. Al final se propone el cambio. Describir antes de juzgar evita acusar a un gráfico que solo estaba mal hecho.}
\milcpaso{3}{milcMagenta}{\textbf{La corrección es pequeña.} Casi siempre basta con una cosa: el eje a cero, quitar el 3D, agrupar tajadas, poner el título. Si tu corrección exige rehacer todo, probablemente no encontraste la trampa principal.}
\milcpaso{4}{milcMagenta}{\textbf{Revisar al otro con las reglas, no con el gusto.} «Me gusta más el azul» no es una revisión. «El eje empieza en 2,5 y hace ver la diferencia más grande de lo que es» sí lo es.}

\begin{drawbox}{milcMagenta}{Antes de entregar}
\checkbox Tres gráficos del mismo conjunto: barras, líneas y circular. \\[1mm] \checkbox Cada uno con título, unidades, fuente y eje vertical desde cero. \\[1mm] \checkbox Un párrafo por gráfico con la pregunta que responde mejor. \\[1mm] \checkbox Un gráfico engañoso real con trampa, impresión falsa y corrección. \\[1mm] \checkbox Revisión de los gráficos de mi pareja con las cinco reglas. \\[1mm] \checkbox Hoja guardada como G8-P1-S7-graficos.xlsx.
\end{drawbox}

\begin{softbox}{milcMagenta}{milcGris}{Plantilla de trabajo}
\textbf{Barras.} \textit{Responde mejor: ¿quién tiene más?} \\[1mm] \textbf{Líneas.} \textit{Responde mejor: ¿cómo cambió en el tiempo?} \\[1mm] \textbf{Circular.} \textit{Responde mejor: ¿qué parte del total?} \\[1mm] \textbf{Gráfico engañoso.} \textit{Trampa, impresión falsa, corrección.} \\[1mm] \textbf{Revisión.} \textit{Regla que le falta al gráfico de mi pareja.}
\end{softbox}

\begin{infoband}{milcMagenta}


\textbf{Lo que va al cuaderno (Actividad 3 · EVALÚA).} \textbf{Título:} «Actividad 3 --- Un gráfico engañoso de verdad». \textbf{Formato:} el gráfico pegado o dibujado, más tres líneas (trampa / impresión falsa / corrección), más la lista de reglas que le faltan a los gráficos de tu pareja. \textbf{Extensión:} media página. \textbf{Sabes que terminaste cuando} nombraste la trampa con una de las seis, tu corrección cabe en una línea, y tu pareja recibió al menos una regla que le faltaba.
\end{infoband}

\puente{Lo que entregas hoy: tres gráficos honestos con su párrafo, y la crítica de uno engañoso. La prueba de calidad: tu pareja lee cada gráfico sin que le expliques nada y dice qué muestra.}

\milcestacion{milcMostaza}{Producto de la sesión}
\textbf{Tres gráficos honestos con su párrafo + crítica de un gráfico engañoso real}.
\begin{softbox}{milcMostaza}{milcGris}{Criterios de calidad}
(1) Los tres gráficos son del mismo conjunto y cumplen las cinco reglas. (2) Cada párrafo nombra una pregunta distinta que ese tipo responde mejor. (3) La crítica nombra una de las seis trampas y propone una corrección concreta. (4) Revisaste los gráficos de tu pareja con las cinco reglas, no con el gusto. (5) La hoja está guardada con nombre claro.
\end{softbox}

\puente{Antes de cerrar, mira lo que hiciste desde cinco lados. Un gráfico con el eje en cero es una forma de respeto por quien lo mira: no le levanta una cerca entre lo que ve y lo que es.}

\milcestacion{milcVino}{Evaluación liberadora · cinco dimensiones}

\begin{softbox}{milcVino}{milcGris}{Sentido de la evaluación}
Evaluar no es solo revisar una nota. Es reconocer qué aprendí, qué cambió en mí, qué emociones manejé, cómo me relaciono con la ciudad y la comunidad, y qué le dejaré a quien viene detrás.
\end{softbox}

\textbf{Mi reflexión liberadora en cinco dimensiones.} Completa cada frase iniciada:

\small
\begin{tabularx}{\textwidth}{>{\bfseries\RaggedRight\arraybackslash}p{3.4cm}Y>{\RaggedRight\arraybackslash}p{4.6cm}}
\rowcolor{milcVino}\color{white}Dimensión & \color{white}\bfseries Mi respuesta (frame starter) & \color{white}\bfseries Algo que descubrí\\
1. Personal & Lo que más cambió en mí fue \linead{6.5cm} & \linead{4.2cm}\\
2. Emocional & La emoción más fuerte fue \linead{2.5cm} y la regulé \linead{3cm} & \linead{4.2cm}\\
3. Ciudadana & En democracia esto importa porque \linead{6cm} & \linead{4.2cm}\\
4. Local (Cartago) & En mi barrio sería útil para \linead{6.5cm} & \linead{4.2cm}\\
5. Intergeneracional & A mi abuela/abuelo le diría \linead{6.5cm} & \linead{4.2cm}\\
\end{tabularx}
\normalsize

\begin{infoband}{milcVino}
\textbf{Para la próxima sustentación.} Vuelve a esta tabla en seis meses. Verás que las respuestas cambian — no porque lo aprendido cambie, sino porque tú cambias. La evaluación liberadora es una conversación contigo a través del tiempo.
\end{infoband}


\puente{Para terminar, tres ideas sobre mostrar sin engañar, contadas con palabras de esta guía: Dussel sobre la información que nunca es neutral, Epicteto sobre describir antes de juzgar, Floridi sobre dar sentido a lo que asusta.}

\milcestacion{milcGradeDark}{Tres ideas para cerrar · Dussel, un estoico y Floridi}

\begin{softbox}{milcVino}{milcGris}{Enrique Dussel · lente del nosotros}
\textit{Es ingenuo creer que la información se lee sola, como si no hubiera conflictos detrás de cada dato.}

{\footnotesize\itshape\color{milcNegro!70}--- idea tomada de Enrique Dussel, contada con palabras de esta guía. No es una cita textual.}

Un gráfico con el eje en 50 no es un error inocente: alguien decidió dónde empezar. En Riosucio, la cerca también la puso alguien. Preguntar quién decidió el eje es leer el gráfico de verdad.

\textbf{¿Quién decidió dónde empieza el eje del gráfico que vi hoy, y qué ganaba con eso?}
\end{softbox}
\linead{\linewidth}\\[1mm]
\linead{\linewidth}

\begin{softbox}{milcMostaza}{milcGris}{Estoicismo · Epicteto · cuidado interior}
\textit{Cuando alguien se baña rápido, no digas que se baña mal: di que se baña rápido. Describe el hecho antes de juzgarlo.}

{\footnotesize\itshape\color{milcNegro!70}--- idea tomada de Epicteto, contada con palabras de esta guía. No es una cita textual.}

Antes de decir «este gráfico engaña», describe: qué tipo es, dónde empieza el eje, cuántas tajadas tiene. La descripción es la que te permite nombrar la trampa sin equivocarte.

\textbf{¿Describí el gráfico antes de juzgarlo, o juzgué primero y describí después?}
\end{softbox}
\linead{\linewidth}\\[1mm]
\linead{\linewidth}

\begin{softbox}{milcTurquesa}{milcGris}{The Onlife Initiative · lente de la infoesfera}
\textit{Tememos y rechazamos aquello a lo que no logramos darle sentido.}

{\footnotesize\itshape\color{milcNegro!70}--- idea tomada de The Onlife Initiative, contada con palabras de esta guía. No es una cita textual.}

Un buen gráfico le da sentido a doscientas filas en un vistazo. Uno malo asusta o confunde, y la gente reacciona al miedo, no al dato. Hacer gráficos honestos es ayudar a otros a entender antes de reaccionar.

\textbf{¿Mi gráfico le da sentido al dato, o solo lo hace ver grande?}
\end{softbox}
\linead{\linewidth}\\[1mm]
\linead{\linewidth}

\begin{infoband}{milcNegro}
\textbf{Nota para el docente.} Las tres ideas están contadas con palabras de esta guía a partir del pensamiento de cada autor; no son citas textuales. De 9.º en adelante se trabajan con la obra y su referencia.
\end{infoband}

\milcestacion{milcVino}{Mi autoevaluación · marca una casilla por fila}

{\small
\setlength{\extrarowheight}{2.2pt}
\begin{tabularx}{\textwidth}{>{\RaggedRight\arraybackslash\bfseries}p{3.0cm}YYY}
\rowcolor{milcVino}\color{white}\bfseries Criterio & \color{white}\bfseries Lo logré & \color{white}\bfseries En proceso & \color{white}\bfseries Necesito apoyo\\[.6mm]
Comprensión del tema & \milcopcion{Explico las ideas con mis palabras.} & \milcopcion{Reconozco las ideas, falta precisión.} & \milcopcion{Necesito repasar lo básico.}\\[.8mm]
\rowcolor{milcGris}Pensamiento computacional & \milcopcion{Usé los pilares al armar mi producto.} & \milcopcion{Usé algunos pilares.} & \milcopcion{Necesito releer la estación 2.}\\[.8mm]
Aplicación y producto & \milcopcion{Mi producto cumple los criterios.} & \milcopcion{Está armado, con ajustes pendientes.} & \milcopcion{Aún está en construcción.}\\[.8mm]
\rowcolor{milcGris}Reflexión 5 dimensiones & \milcopcion{Respondí cada una con honestidad.} & \milcopcion{Respondí algunas con detalle.} & \milcopcion{Mis respuestas son muy generales.}\\[.8mm]
Triángulo de pensamiento & \milcopcion{Conecté las 3 voces con mi trabajo.} & \milcopcion{Conecté al menos una.} & \milcopcion{No las relaciono con mi proceso.}\\[.8mm]
\end{tabularx}}

\vspace{3mm}
\textbf{Hoy entendí que:}\\[1mm]
\linead{\linewidth}\\[1mm]
\linead{\linewidth}

\textbf{Mi compromiso para la próxima sesión es:}\\[1mm]
\linead{\linewidth}\\[1mm]
\linead{\linewidth}

\begin{softbox}{milcMostaza}{milcGris}{Cita de despedida · Marco Aurelio}
\textit{``El obstáculo en el camino se vuelve el camino. Lo que estorba al camino se vuelve camino.''}\\[1mm]
Cada error de hoy, bien observado, se convierte en pieza del camino que pisarás mañana.
\end{softbox}

\small
\begin{tabularx}{\textwidth}{>{\bfseries}p{3.6cm}Y}
\rowcolor{milcGradeDark!12}Nombre completo: & \linead{10cm}\\
Fecha: & \linead{4cm} \quad Curso/grupo: \linead{4cm}\\
Mi firma: & \linead{9cm}\\
\end{tabularx}
\normalsize

\begin{infoband}{milcGradeDark}
\textbf{Para la próxima sesión.} Llega con tus tres gráficos y la crítica del engañoso. En la sesión 8 vas a hacer que la hoja avise sola con colores cuando un valor se sale de rango. Y que rechace los datos que no deben entrar.
\end{infoband}

\end{document}
